\section{An exact third-height stratum in characteristic two}
A constant Boolean formula provides the next nonzero height quotient.
\begin{proposition}[An exact height-three binary gadget]
Let
\[
 \Psi_3=(\neg x_1\vee\neg x_2)\wedge
 (\neg x_1\vee\neg x_3)\wedge
 (\neg x_1\vee\neg x_2\vee\neg x_3)\wedge
 (x_1\vee\neg x_2\vee\neg x_3).
\]
Then $S_{\Psi_3}=4$ and its incidence replica coefficients are
\[
 T_3=1002460\equiv0\pmod4,\qquad
 T_7=1382535756736336500\equiv4\pmod8.
\]
Its counting hypersurface in characteristic two has height exactly
three.  It has ten incidences, $u=21$, and original stabilization
parameter $Q=32$.
\end{proposition}
\begin{proof}
The Boolean function is ``at most one of $x_1,x_2,x_3$ is true'',
so its satisfying count is four.  Here is a finite integer formula
and evaluation certificate for the two further coefficients.
Set $\binom nk=0$ outside $0\leq k\leq n$, and put
\[
 A_M(a,b)=
 \begin{cases}
 \dfrac{M!}{a!b!(M-a-b)!},&a+b\leq M,\\
 0,&a+b>M,
 \end{cases}
\]
and, for $f\in\{0,1\}^3$, put
\[
 C_{M,f}(a,b,c)=\sum_{j=0}^M(-1)^j\binom Mj
 \binom{M-j}{a-jf_1}\binom{M-j}{b-jf_2}
 \binom{M-j}{c-jf_3}.
\]
The pair clause permits exactly the three local assignments
$00,10,01$, so its $M$-replica coefficient is $A_M$.
A ternary clause with unique forbidden assignment $f$ has local
polynomial
$(1+z_1)(1+z_2)(1+z_3)-z_1^{f_1}z_2^{f_2}z_3^{f_3}$.
The binomial expansion of its $M$th power gives $C_{M,f}$.
Multiplying the two pair clauses, the two ternary clauses, and
the variable marginal counts therefore gives the exact formula
\[
 T_M=\sum_{a,b,c=0}^M
 \binom Ma\binom Mb\binom Mc\,
 A_M(a,b)A_M(a,c)C_{M,(1,1,1)}(a,b,c)
 C_{M,(0,1,1)}(a,b,c).
\]
For transparency, grouping this finite sum by $a$ gives the
following integer certificate; the absent $M=3$ entries are zero.
\[
\begin{array}{c|r|r}
 a&M=3&M=7\\\hline
 0&9075&665953307379\\
 1&857547&5184197281860583\\
 2&135837&377589783130120323\\
 3&1&909264005602246875\\
 4&0&90276840331103125\\
 5&0&220262726359053\\
 6&0&1711339161\\
 7&0&1\\\hline
 \text{sum}&1002460&1382535756736336500
\end{array}
\]
Summing the displayed entries gives the asserted residues. The early-stabilization
formula gives $a_2=-S$, $a_4=-T_3$, and $a_8=-T_7$.
The first two height tests vanish and $a_8/4$ is odd, proving the claim.
\end{proof}

\begin{lemma}[The diagonal replica residue]
For the incidence polynomial of any input formula and every prime $p$,
\[
 T_{1+p+\cdots+p^{k-1}}\equiv S^k\pmod p.
\]
\end{lemma}
\begin{proof}
Modulo $p$, the relevant power of $P$ is
$\prod_{j=0}^{k-1}P(z^{p^j})$.  Every individual incidence variable
has exponent at most two in a monomial of $P$.  To obtain target
exponent $1+p+\cdots+p^{k-1}$, reduction modulo $p$ forces the
lowest exponent to be one; this holds also for $p=2$, since the
possibilities are $0,1,2$.  There is no carry.  Induction on the
digits forces exponent one in every factor and every variable.
The target coefficient is therefore the product of $k$ copies of
$[Z]P=S$, modulo $p$.
\end{proof}

\begin{corollary}[The next binary height decision]
On the promise that the counting hypersurface has height at least
three, deciding height three versus height at least four, including
infinity, is $\oplus\mathrm P$-complete.
\end{corollary}
\begin{proof}
For hardness conjoin an arbitrary input $\Phi$ with $\Psi_3$ on
disjoint variables.  Its satisfying count is even and its $T_3$ is
divisible by four, so its height is at least three.  Its height is
three precisely when
\[
 T_7(\Phi)T_7(\Psi_3)/4\equiv T_7(\Phi)
 \equiv S_\Phi\pmod2
\]
is nonzero.  All three early coefficient identities apply to the
combined construction.  For membership $T_7$ counts polynomial-size
configurations, since seven is fixed.  Guess four distinct such
configurations in increasing lexicographic order.  The number of
accepting guesses is $\binom{T_7}{4}$, whose parity is the binary
digit of $T_7$ in position two.  On the promise, this is exactly
$T_7/4\bmod2$, the third height test.  Thus it is a \#P parity.
\end{proof}

